
\chapter{Related Work}
\label{chap:related-work}

The development of a generative surrogate for Rayleigh--Taylor instability must be understood against two constraints that pull in opposite directions. The surrogate must remain faithful to the multi-scale structure of the flow, but it must also generate realizations substantially faster than a high-fidelity simulation. This chapter reviews the scientific and methodological context of the present work as a continuous argument: DNS provides the reference distribution, multi-resolution representations offer a way to expose its structure, and score-based models provide a mechanism for sampling it.

Direct Numerical Simulation is a reference approach in turbulence research because it resolves the relevant spatial and temporal scales without introducing an explicit turbulence closure model \citep{MoinMahesh1998}. For RTI, this makes it possible to follow the evolution from initial interface perturbations to non-linear bubbles, spikes, and turbulent mixing structures. Reviews of RTI emphasize the importance of this instability for the development of flow, turbulence, and mixing across a broad range of regimes \citep{KULL1991197,zhou2017rti1}. DNS is therefore well suited to produce the training and validation fields required by this study. Its computational cost, which grows with spatial resolution, simulated time, and physical-model complexity, limits the number of independent realizations that can be generated. This limitation is central here: a small DNS ensemble may contain rich physics, but it is not directly compatible with applications that require a large number of samples.

This cost motivates the use of reduced-order and data-driven surrogates in computational science. Multifidelity and surrogate methods are commonly introduced precisely because uncertainty propagation, inference, optimization, and related outer-loop problems require many model evaluations \citep{Peherstorfer2018}. In fluid mechanics, machine learning has been used for model reduction, reconstruction, forecasting, control, and the extraction of coherent structures, while its limitations in extrapolation, physical consistency, and data requirements remain important \citep{Brunton2020MLFluid}. The present work addresses a specific part of this broader problem: learning a generative model of RTI density fields so that the surrogate can produce diverse realizations for statistical analysis, rather than only predicting a single state from a prescribed input.

The representation used by the surrogate is particularly important because RTI fields contain structures over a range of spatial scales. Fourier decompositions are natural for periodic domains and provide a global description in terms of spatial frequencies. They are especially useful for spectral diagnostics and for identifying how fluctuations are distributed across scales. However, each Fourier mode is spatially global, which makes the representation less localized when the relevant structures are sharp, intermittent, or confined to a limited region. Wavelet analysis provides a complementary description localized in both position and scale. Its use in turbulence analysis is well established \citep{Farge1992}, and the mathematical foundations of multi-resolution analysis are developed in the work of Mallat and related references \citep{Mallat2009,Meyer1992}. For RTI, this distinction is meaningful because the large-scale deformation of the mixing interface coexists with localized gradients and fine-scale structures.

Score-based generative models offer a probabilistic framework adapted to this representation-learning problem. In the continuous-time formulation of Song et al., a forward stochastic differential equation progressively perturbs the data distribution toward a Gaussian prior, while a neural network estimates the score of the perturbed distribution \citep{Song2021SDE}. Sampling reverses this process through a time-reversed SDE or through the associated probability-flow ODE. The computational cost of sampling is proportional to the number of time discretization steps and, consequently, to the number of score-network evaluations. This creates a tension between accuracy and throughput: coarse integration may be fast but inaccurate, whereas fine integration improves the approximation of the reverse dynamics at a higher cost. The mathematical basis for time reversal is given by Anderson's formulation of reverse-time diffusion processes \citep{ANDERSON1982313}.

The WSGM of Guth et al. directly addresses this tension by applying score-based models to normalized wavelet coefficients conditioned across scales \citep{Guth2022WSGM}. Instead of learning one global reverse diffusion for the full-resolution image, WSGM first samples the coarsest representation and then generates finer detail coefficients conditionally. The wavelet normalization reduces the ill-conditioning associated with the score at fine resolution. For Gaussian multiscale distributions, the authors prove that the discretization schedule needed to reach a fixed error does not have to grow with image size. Their experiments on a physical process at phase transition and on natural images show that the WSGM reaches comparable quality with far fewer steps than a conventional SGM, with an acceleration of up to approximately two orders of magnitude in the reported experiments. The result is significant for the present application because it links the multi-scale structure of the data to the computational cost of generation. It must nevertheless be interpreted within its scope: the theoretical result is not a general convergence theorem for arbitrary non-Gaussian turbulent fields, and the quality of the adaptation depends on the statistics of the RTI wavelet coefficients.

The approach developed in this report follows this coarse-to-fine rationale while adapting it to the available RTI data and implementation constraints. The direct image-space SGM establishes a necessary baseline, making it possible to separate the benefit of the representation from the benefit of the diffusion model itself. The wavelet pipeline decomposes each field into an approximation component and directional detail components, normalizes the coefficient channels, and learns a conditional score model for the details given the coarse information. The reconstructed fields are then compared with DNS references using both generic feature-space metrics and physical diagnostics. The contribution of this study is consequently not to assume that WSGM transfers unchanged to RTI, but to test whether its mechanism provides a useful basis for a fast and statistically credible generative surrogate of a multi-scale hydrodynamic instability.
